\documentclass[11pt,letterpaper]{article}

% ------------------------------------------------------------
% Paquetes
% ------------------------------------------------------------
\usepackage[utf8]{inputenc}
\usepackage[T1]{fontenc}
\usepackage[spanish,es-tabla,es-noshorthands,es-nodecimaldot]{babel}
\usepackage{lmodern}
\usepackage[margin=2.5cm]{geometry}
\usepackage{amsmath,amssymb,amsthm}
\usepackage{graphicx}
\usepackage{float}
\usepackage{booktabs}
\usepackage{array}
\usepackage{multirow}
\usepackage{xcolor}
\usepackage{listings}
\usepackage{algorithm}
\usepackage{algpseudocode}
\usepackage{enumitem}
\usepackage{fancyhdr}
\usepackage[hidelinks]{hyperref}

% ------------------------------------------------------------
% Configuración general
% ------------------------------------------------------------
\setlength{\parskip}{0.5em}
\setlength{\parindent}{0pt}

\pagestyle{fancy}
\fancyhf{}
\lhead{Análisis y Diseño de Algoritmos II}
\rhead{Proyecto ModEx}
\cfoot{\thepage}

\floatname{algorithm}{Algoritmo}
\renewcommand{\listalgorithmname}{Lista de algoritmos}
\algrenewcommand\algorithmicrequire{\textbf{Entrada:}}
\algrenewcommand\algorithmicensure{\textbf{Salida:}}

\lstset{
  language=Python,
  basicstyle=\ttfamily\small,
  keywordstyle=\color{blue!70!black}\bfseries,
  commentstyle=\color{green!50!black}\itshape,
  stringstyle=\color{red!60!black},
  numbers=left,
  numberstyle=\tiny\color{gray},
  frame=single,
  breaklines=true,
  showstringspaces=false,
  tabsize=4,
  literate={á}{{\'a}}1 {é}{{\'e}}1 {í}{{\'i}}1 {ó}{{\'o}}1 {ú}{{\'u}}1 {ñ}{{\~n}}1
}

% Marcador para las secciones pendientes de redactar
\newcommand{\introducirtexto}{%
  \par\noindent\fcolorbox{red!70!black}{red!5}{%
    \parbox{\dimexpr\linewidth-2\fboxsep-2\fboxrule}{%
      \centering\vspace{0.4em}\textbf{\color{red!70!black}INTRODUCIR TEXTO}\vspace{0.4em}}}\par}

% Notación del problema
\newcommand{\RS}{\mathit{RS}}
\newcommand{\Rmax}{R_{\max}}
\newcommand{\Ext}{\operatorname{Ext}}
\newcommand{\Esf}{\operatorname{Esfuerzo}}
\newcommand{\Mod}{\operatorname{Mod}}

% ------------------------------------------------------------
% Documento
% ------------------------------------------------------------
\begin{document}

% ------------------------------------------------------------
% Portada
% ------------------------------------------------------------
\begin{titlepage}
  \centering
  \vspace*{1cm}
  {\large Universidad del Valle\par}
  {\large Escuela de Ingeniería de Sistemas y Computación\par}
  \vspace{0.5cm}
  {\large Análisis y Diseño de Algoritmos II\par}
  \vspace{3cm}
  {\Huge\bfseries Moderando el extremismo de opiniones en una red social\par}
  \vspace{0.8cm}
  {\Large Proyecto de curso -- Informe\par}
  \vspace{3cm}
  {\large\textbf{Integrantes}\par}
  \vspace{0.3cm}
  \begin{tabular}{ll}
    Santiago Arias Rojas        & 2416285 \\
    Sebastián Calvo Carvajal    & \rule{2cm}{0.4pt} \\
    Sebastián Rubio Piedrahita  & \rule{2cm}{0.4pt} \\
  \end{tabular}
  \vfill
  {\large\textbf{Profesor:} Jesús Alexander Aranda\par}
  {\large\textbf{Monitor:} Samuel Galindo\par}
  \vspace{1cm}
  {\large Noviembre de 2026\par}
\end{titlepage}

% ------------------------------------------------------------
% Tabla de contenido
% ------------------------------------------------------------
\tableofcontents
\newpage
\listoftables
\listofalgorithms
\newpage

% ============================================================
\section{Introducción}
% ============================================================

El presente informe documenta el desarrollo del proyecto de curso de
Análisis y Diseño de Algoritmos II, cuyo objetivo es aplicar y comparar tres
estrategias de diseño de algoritmos ---fuerza bruta, algoritmos voraces y
programación dinámica--- sobre un mismo problema de naturaleza combinatoria:
el problema de \emph{moderar el extremismo de opiniones en una red social}
(ModEx).

Para cada estrategia se describe el algoritmo diseñado, se analiza su
complejidad temporal (y espacial cuando aplica) y se argumenta si conduce o no
a una solución óptima. Finalmente se comparan experimentalmente las tres
soluciones implementadas y se concluye sobre las ventajas y desventajas de
utilizar cada enfoque en la práctica.

El documento se organiza así: la Sección~\ref{sec:problema} presenta y
formaliza el problema; las Secciones~\ref{sec:fb}, \ref{sec:voraz} y
\ref{sec:pd} desarrollan las soluciones por fuerza bruta, voraz y programación
dinámica, respectivamente; la Sección~\ref{sec:implementacion} describe la
implementación; la Sección~\ref{sec:comparacion} compara los resultados
obtenidos y la Sección~\ref{sec:conclusiones} presenta las conclusiones.

% ============================================================
\section{El problema ModEx}
\label{sec:problema}
% ============================================================

\subsection{Descripción}

Se considera una red social formada por $n$ agentes, cada uno con una opinión
respecto a una afirmación sobre un tema. La opinión se cuantifica como un
entero entre $-100$ (totalmente desfavorable) y $100$ (completamente
favorable); entre más cercana a $0$ esté la opinión, más moderada es, y entre
más alejada de $0$, más extrema.

Cada agente posee además:
\begin{itemize}
  \item Un nivel de \textbf{receptividad} $r_i \in [0,1]$, que modela qué tan
        fácil es moderar su opinión ($1$: muy fácil, $0$: muy difícil).
  \item Un nivel de \textbf{influencia} $\mathit{inf}_i \in [1,3]$, que modela el
        impacto de su opinión en la red.
\end{itemize}

Moderar la opinión de un agente requiere invertir un esfuerzo que depende de
su opinión, su receptividad y su influencia. La red cuenta con un valor
máximo $\Rmax$ de esfuerzo disponible. El problema consiste en decidir sobre
cuáles agentes concentrar los esfuerzos de moderación para reducir lo más
posible el extremismo de la red sin superar $\Rmax$.

\subsection{Formalización}

Una red social $\RS$ es una pareja $\langle AG, \Rmax \rangle$, donde
$AG = \langle a_0, \dots, a_{n-1} \rangle$ es una secuencia de agentes y
$\Rmax \in \mathbb{N}$ es el valor entero máximo disponible para moderar.
Cada agente es una terna
$a_i = \langle o_i, r_i, \mathit{inf}_i \rangle$ con
$o_i \in \{-100,\dots,100\}$, $r_i \in [0,1]$ y $1 \le \mathit{inf}_i \le 3$.

\paragraph{Extremismo.} El extremismo de una red se define como
\begin{equation}
  \Ext(\RS) = \frac{\sqrt{\sum_{i=0}^{n-1} o_i^2 \cdot \mathit{inf}_i}}{n}.
  \label{eq:ext}
\end{equation}

\paragraph{Estrategia de moderación.} Es una secuencia
$E = \langle e_0, \dots, e_{n-1} \rangle$ con $e_i \in \{0,1\}$, donde
$e_i = 1$ indica que la opinión del agente $a_i$ se modera y $e_i = 0$ que no.

\paragraph{Aplicación de una estrategia.} $\Mod(\RS, E) = \RS'$ es la red con
la misma estructura de $\RS$ en la que
\[
  o'_i =
  \begin{cases}
    o_i & \text{si } e_i = 0,\\
    0   & \text{si } e_i = 1.
  \end{cases}
\]

\paragraph{Esfuerzo.} El esfuerzo de moderar $\RS$ con la estrategia $E$ es
\begin{equation}
  \Esf(\RS, E) = \sum_{i=0}^{n-1}
  \left\lceil |o_i - o'_i| \cdot (1 - r_i) \cdot \mathit{inf}_i \right\rceil .
  \label{eq:esf}
\end{equation}
Una estrategia es \emph{aplicable} si y solo si $\Esf(\RS, E) \le \Rmax$.

\paragraph{Problema ModEx.}
\begin{description}[leftmargin=2cm]
  \item[Entrada:] una red social $\RS = \langle AG, \Rmax \rangle$.
  \item[Salida:] una estrategia de moderación $E$ aplicable tal que
        $\Ext(\Mod(\RS, E))$ sea mínimo.
\end{description}

\subsection{Ejemplos}
\label{sec:ejemplos}

\subsubsection*{Ejemplo 1}
$\RS_1 = (\langle\langle 100, 0.5, 1\rangle, \langle 100, 0.1, 1\rangle,
\langle -10, 0.1, 1\rangle\rangle, 55)$.

\begin{itemize}
  \item $E_1 = \langle 0,0,1 \rangle$:
        $\Esf = \lceil 10 \cdot 0.9 \cdot 1 \rceil = 9 \le 55$ (aplicable) y
        $\Ext = \sqrt{100^2 + 100^2 + 0^2}/3 \approx 47.14$.
  \item $E_2 = \langle 1,0,0 \rangle$:
        $\Esf = \lceil 100 \cdot 0.5 \cdot 1 \rceil = 50 \le 55$ (aplicable) y
        $\Ext = \sqrt{0^2 + 100^2 + (-10)^2}/3 \approx 33.50$.
\end{itemize}

\subsubsection*{Ejemplo 2}
$\RS_2 = (\langle\langle -30, 0.9, 1\rangle, \langle 20, 0.1, 2\rangle,
\langle 50, 0.5, 1\rangle\rangle, 35)$.

\begin{itemize}
  \item $E_1 = \langle 0,1,0 \rangle$:
        $\Esf = \lceil 20 \cdot 0.9 \cdot 2 \rceil = 36 > 35$, por lo que
        \emph{no} es aplicable.
  \item $E_2 = \langle 1,0,1 \rangle$:
        $\Esf = \lceil 30 \cdot 0.1 \cdot 1 \rceil + \lceil 50 \cdot 0.5 \cdot 1 \rceil
        = 28 \le 35$ (aplicable) y
        $\Ext = \sqrt{0^2 + 20^2 \cdot 2 + 0^2}/3 \approx 9.43$.
\end{itemize}

% ============================================================
\section{Solución por fuerza bruta}
\label{sec:fb}
% ============================================================

\subsection{Descripción del algoritmo}

El algoritmo de fuerza bruta genera todas las estrategias de moderación
posibles ---es decir, todas las secuencias binarias de longitud $n$---, calcula
para cada una su esfuerzo y el extremismo de la red moderada, descarta las que
no son aplicables y, entre las aplicables, retorna la de menor extremismo.

La generación de estrategias se hace recursivamente: las estrategias de
longitud $n$ se obtienen anteponiendo un $1$ y un $0$ a cada una de las
estrategias de longitud $n-1$, con caso base la secuencia vacía.

\begin{algorithm}[H]
  \caption{Generación de todas las estrategias}
  \label{alg:generar}
  \begin{algorithmic}[1]
    \Function{GenerarEstrategias}{$n$}
      \If{$n \le 0$}
        \State \Return $[\,[\,]\,]$
      \EndIf
      \State $S \gets$ \Call{GenerarEstrategias}{$n-1$}
      \State \Return $\{\, [1] + s : s \in S \,\} \cup \{\, [0] + s : s \in S \,\}$
    \EndFunction
  \end{algorithmic}
\end{algorithm}

\begin{algorithm}[H]
  \caption{ModEx por fuerza bruta}
  \label{alg:fb}
  \begin{algorithmic}[1]
    \Require Red social $\RS = \langle AG, \Rmax \rangle$ con $n$ agentes
    \Ensure $(E^*, \Esf(\RS,E^*), \Ext(\Mod(\RS,E^*)))$
    \State $E^* \gets \textsc{nil}$;\quad $\mathit{mejorExt} \gets \infty$;\quad $\mathit{mejorEsf} \gets 0$
    \ForAll{$E \in$ \Call{GenerarEstrategias}{$n$}}
      \State $\mathit{esf} \gets \Esf(\RS, E)$
      \State $\mathit{ext} \gets \Ext(\Mod(\RS, E))$
      \If{$\mathit{esf} \le \Rmax$ \textbf{and} $\mathit{ext} < \mathit{mejorExt}$}
        \State $E^* \gets E$;\quad $\mathit{mejorExt} \gets \mathit{ext}$;\quad $\mathit{mejorEsf} \gets \mathit{esf}$
      \EndIf
    \EndFor
    \State \Return $(E^*, \mathit{mejorEsf}, \mathit{mejorExt})$
  \end{algorithmic}
\end{algorithm}

\subsection{Complejidad}

\paragraph{Generación de estrategias.} Sea $G(n)$ el costo de
\textsc{GenerarEstrategias}$(n)$. La llamada recursiva retorna $2^{n-1}$
listas de longitud $n-1$, y construir las $2^n$ nuevas listas de longitud $n$
cuesta $\Theta(n \cdot 2^n)$. Así,
\[
  G(n) = G(n-1) + \Theta(n \cdot 2^n)
  \quad\Longrightarrow\quad
  G(n) = \sum_{k=1}^{n} \Theta(k \cdot 2^k) = \Theta(n \cdot 2^n).
\]

\paragraph{Evaluación.} Hay $2^n$ estrategias, y para cada una el cálculo del
esfuerzo (Ecuación~\ref{eq:esf}) y del extremismo (Ecuación~\ref{eq:ext})
recorre los $n$ agentes realizando operaciones de costo constante, es decir,
$\Theta(n)$ por estrategia y $\Theta(n \cdot 2^n)$ en total.

\paragraph{Total.} La complejidad temporal del algoritmo es
\[
  T(n) = \Theta(n \cdot 2^n).
\]

\paragraph{Espacio.} Como la implementación materializa todas las estrategias
en memoria antes de evaluarlas, la complejidad espacial también es
$\Theta(n \cdot 2^n)$. Esta podría reducirse a $\Theta(n)$ generando las
estrategias una a una (por ejemplo, con un generador o recorriendo los
enteros de $0$ a $2^n - 1$ en binario), sin alterar la complejidad temporal.

\subsection{Corrección}

El algoritmo siempre da la respuesta correcta al problema:

\begin{enumerate}
  \item \textbf{Exhaustividad.} El espacio de soluciones de ModEx es el
        conjunto $\{0,1\}^n$ de todas las estrategias. Por inducción sobre $n$,
        \textsc{GenerarEstrategias}$(n)$ produce exactamente las $2^n$
        secuencias binarias de longitud $n$: para $n=0$ produce la secuencia
        vacía, y si para $n-1$ produce todas las secuencias de longitud $n-1$,
        al anteponer $0$ y $1$ a cada una se obtienen todas las de longitud $n$,
        sin repeticiones.
  \item \textbf{Factibilidad.} Solo se acepta una estrategia como candidata si
        $\Esf(\RS,E) \le \Rmax$, por lo que la estrategia retornada siempre es
        aplicable.
  \item \textbf{Existencia.} La estrategia $\langle 0,\dots,0 \rangle$ tiene
        esfuerzo $0 \le \Rmax$, de modo que siempre existe al menos una
        estrategia aplicable y el algoritmo nunca retorna \textsc{nil}.
  \item \textbf{Optimalidad.} Se mantiene el invariante de que, tras procesar
        un subconjunto de estrategias, $E^*$ es la de menor extremismo entre las
        aplicables procesadas. Al terminar el recorrido se han procesado todas,
        luego $E^*$ minimiza $\Ext(\Mod(\RS,E))$ sobre todas las estrategias
        aplicables.
\end{enumerate}

En caso de empate en el extremismo, el algoritmo conserva la primera
estrategia encontrada, lo cual no afecta la optimalidad.

\subsection{Resultados sobre los ejemplos}

La Tabla~\ref{tab:fb-ejemplos} muestra la salida del algoritmo de fuerza bruta
para los ejemplos de la Sección~\ref{sec:ejemplos}. Al ser exhaustivo, estos
valores corresponden a las soluciones óptimas y sirven como referencia para
evaluar las demás estrategias.

\begin{table}[H]
  \centering
  \caption{Salida del algoritmo de fuerza bruta sobre los ejemplos del enunciado}
  \label{tab:fb-ejemplos}
  \begin{tabular}{lccc}
    \toprule
    Entrada & $E^*$ & $\Esf(\RS,E^*)$ & $\Ext(\Mod(\RS,E^*))$ \\
    \midrule
    $\RS_1$ & $\langle 1,0,0 \rangle$ & 50 & 33.50 \\
    $\RS_2$ & $\langle 1,0,1 \rangle$ & 28 & 9.43 \\
    \bottomrule
  \end{tabular}
\end{table}

% ============================================================
\section{Solución voraz}
\label{sec:voraz}
% ============================================================

\subsection{Descripción del algoritmo}
\introducirtexto

\subsection{Entendimos el algoritmo}
\label{sec:voraz-entendimos}

\subsubsection{Ejemplos del enunciado}
\introducirtexto

\subsubsection{Entradas adicionales}
\introducirtexto

\begin{table}[H]
  \centering
  \caption{Resultados del algoritmo voraz frente a la solución óptima}
  \label{tab:voraz-entradas}
  \begin{tabular}{lcccccc}
    \toprule
    \multirow{2}{*}{Entrada} & \multicolumn{3}{c}{Voraz} & \multicolumn{2}{c}{Óptimo} & \multirow{2}{*}{¿Óptima?} \\
    \cmidrule(lr){2-4}\cmidrule(lr){5-6}
     & $E$ & Esf. & Ext. & Esf. & Ext. & \\
    \midrule
    $\RS_1$ & -- & -- & -- & 50 & 33.50 & -- \\
    $\RS_2$ & -- & -- & -- & 28 & 9.43  & -- \\
    $\RS_3$ & -- & -- & -- & -- & --    & -- \\
    $\RS_4$ & -- & -- & -- & -- & --    & -- \\
    $\RS_5$ & -- & -- & -- & -- & --    & -- \\
    $\RS_6$ & -- & -- & -- & -- & --    & -- \\
    \bottomrule
  \end{tabular}
\end{table}

\subsection{Complejidad}
\introducirtexto

\subsection{Corrección}
\introducirtexto

% ============================================================
\section{Solución por programación dinámica}
\label{sec:pd}
% ============================================================

\subsection{Caracterización de la estructura de una solución óptima}
\introducirtexto

\subsection{Definición recursiva del valor de una solución óptima}
\introducirtexto

\subsection{Algoritmo para calcular el costo de una solución óptima}
\introducirtexto

\subsection{Algoritmo para calcular una solución óptima}
\introducirtexto

\subsection{Complejidad}

\subsubsection{Complejidad en tiempo}
\introducirtexto

\subsubsection{Complejidad en espacio}
\introducirtexto

\subsubsection{¿Es útil en la práctica?}
\introducirtexto

% ============================================================
\section{Implementación}
\label{sec:implementacion}
% ============================================================

\subsection{Lenguaje y estructura del proyecto}

El proyecto se implementó en Python~3. El código fuente se organiza así:

\begin{itemize}
  \item \texttt{agente.py}: clase \texttt{Agente}, que encapsula la opinión,
        la receptividad y la influencia de un agente. Si no se suministran
        valores, se generan aleatoriamente dentro de los rangos del problema
        (útil para generar casos de prueba).
  \item \texttt{red.py}: clase \texttt{Red}, que representa la red social
        $\langle AG, \Rmax \rangle$ y calcula el extremismo de la red,
        opcionalmente tras aplicar una estrategia de moderación.
  \item \texttt{estrategias/fuerzabruta.py}: solución por fuerza bruta
        (cálculo del esfuerzo, generación de estrategias y búsqueda
        exhaustiva).
  \item \texttt{estrategias/voraz.py}: solución voraz.
  \item \texttt{estrategias/dinamica.py}: solución por programación dinámica.
  \item \texttt{pruebas.py}: casos de prueba.
  \item \texttt{main.py}: punto de entrada del programa.
\end{itemize}

\subsection{Formato de entrada y salida}

\paragraph{Entrada.} Archivo de texto con $n+2$ líneas: la primera contiene
$n$; las siguientes $n$ contienen \texttt{op,recep,inf} de cada agente; la
última contiene $\Rmax$.

\begin{lstlisting}[language={},numbers=none,caption={Formato del archivo de entrada}]
n
op0,recep0,inf0
...
op(n-1),recep(n-1),inf(n-1)
R_max
\end{lstlisting}

\paragraph{Salida.} Archivo de texto con $n+2$ líneas: la primera contiene el
extremismo de la red moderada, la segunda el esfuerzo de la estrategia y las
$n$ restantes indican si cada agente fue moderado ($1$) o no ($0$).

\begin{lstlisting}[language={},numbers=none,caption={Formato del archivo de salida}]
Ext
Esf
mod0
...
mod(n-1)
\end{lstlisting}

\subsection{Interfaz}
\introducirtexto

% ============================================================
\section{Comparación de resultados}
\label{sec:comparacion}
% ============================================================

\subsection{Diseño de los casos de prueba}

Los casos de prueba se diseñaron para cubrir tanto situaciones típicas como
casos borde del problema:

\begin{itemize}
  \item \textbf{Ejemplos del enunciado} ($\RS_1$ y $\RS_2$), cuyos resultados
        pueden verificarse a mano.
  \item \textbf{Presupuesto nulo} ($\Rmax = 0$): solo pueden moderarse
        agentes cuyo esfuerzo sea $0$.
  \item \textbf{Receptividad máxima} ($r_i = 1$): moderar al agente no cuesta
        esfuerzo, por lo que debería moderarse siempre.
  \item \textbf{Receptividad mínima} ($r_i = 0$): el esfuerzo es máximo para
        la opinión e influencia del agente.
  \item \textbf{Presupuesto holgado}: $\Rmax$ suficiente para moderar a
        todos los agentes, en cuyo caso el extremismo óptimo es $0$.
  \item \textbf{Presupuesto ajustado}: $\Rmax$ igual o una unidad por debajo
        del esfuerzo de una estrategia candidata, para verificar el manejo de
        la restricción $\Esf(\RS,E) \le \Rmax$ y del redondeo hacia arriba.
\end{itemize}

La Tabla~\ref{tab:casos} resume los casos diseñados.

\begin{table}[H]
  \centering
  \caption{Casos de prueba diseñados}
  \label{tab:casos}
  \begin{tabular}{clc>{\raggedright\arraybackslash}p{5cm}}
    \toprule
    Caso & Agentes $\langle o_i, r_i, \mathit{inf}_i \rangle$ & $\Rmax$ & Propósito \\
    \midrule
    1 & $\langle 100,0.5,1\rangle, \langle 100,0.1,1\rangle, \langle -10,0.1,1\rangle$ & 55  & Ejemplo 1 del enunciado \\
    2 & $\langle -30,0.9,1\rangle, \langle 20,0.1,2\rangle, \langle 50,0.5,1\rangle$   & 35  & Ejemplo 2 del enunciado \\
    3 & $\langle 50,0.5,1\rangle, \langle -40,0.2,2\rangle, \langle 10,0.9,3\rangle$   & 0   & Presupuesto nulo \\
    4 & $\langle 100,1,3\rangle, \langle -100,1,3\rangle$                              & 0   & Receptividad máxima \\
    5 & $\langle 40,0.5,2\rangle, \langle 30,0,1\rangle$                               & 40  & Presupuesto exacto \\
    6 & $\langle 90,0,1\rangle, \langle -60,0.5,3\rangle$                              & 90  & Receptividad mínima \\
    7 & $\langle 7,0.3,1.5\rangle, \langle 5,0.8,1\rangle$                             & 8   & Redondeo del esfuerzo \\
    8 & $\langle 7,0.3,1.5\rangle, \langle 5,0.8,1\rangle$                             & 7   & Presupuesto una unidad menor \\
    9 & $\langle -100,0.2,3\rangle, \langle 80,0.6,2\rangle, \langle -20,0,1.2\rangle$ & 500 & Presupuesto holgado \\
    \bottomrule
  \end{tabular}
\end{table}

\subsection{Resultados}

La Tabla~\ref{tab:comparacion} compara el extremismo obtenido por cada
estrategia en los casos de prueba. Los valores de fuerza bruta corresponden al
óptimo.

\begin{table}[H]
  \centering
  \caption{Comparación de las soluciones obtenidas por cada estrategia}
  \label{tab:comparacion}
  \begin{tabular}{cccccccc}
    \toprule
    \multirow{2}{*}{Caso} & \multicolumn{3}{c}{Fuerza bruta} & \multicolumn{2}{c}{Voraz} & \multicolumn{2}{c}{Prog. dinámica} \\
    \cmidrule(lr){2-4}\cmidrule(lr){5-6}\cmidrule(lr){7-8}
     & $E$ & Esf. & Ext. & Esf. & Ext. & Esf. & Ext. \\
    \midrule
    1 & $\langle 1,0,0\rangle$ & 50  & 33.50 & -- & -- & -- & -- \\
    2 & $\langle 1,0,1\rangle$ & 28  & 9.43  & -- & -- & -- & -- \\
    3 & $\langle 0,0,0\rangle$ & 0   & 25.82 & -- & -- & -- & -- \\
    4 & $\langle 1,1\rangle$   & 0   & 0.00  & -- & -- & -- & -- \\
    5 & $\langle 1,0\rangle$   & 40  & 15.00 & -- & -- & -- & -- \\
    6 & $\langle 0,1\rangle$   & 90  & 45.00 & -- & -- & -- & -- \\
    7 & $\langle 1,0\rangle$   & 8   & 2.50  & -- & -- & -- & -- \\
    8 & $\langle 0,1\rangle$   & 1   & 4.29  & -- & -- & -- & -- \\
    9 & $\langle 1,1,1\rangle$ & 328 & 0.00  & -- & -- & -- & -- \\
    \bottomrule
  \end{tabular}
\end{table}

\subsection{Tiempos de ejecución}

La Tabla~\ref{tab:tiempos} muestra el tiempo de ejecución de cada estrategia
para redes generadas aleatoriamente de distintos tamaños.

\begin{table}[H]
  \centering
  \caption{Tiempo de ejecución (en segundos) según el número de agentes}
  \label{tab:tiempos}
  \begin{tabular}{cccc}
    \toprule
    $n$ & Fuerza bruta & Voraz & Prog. dinámica \\
    \midrule
    5   & -- & -- & -- \\
    10  & -- & -- & -- \\
    15  & -- & -- & -- \\
    20  & -- & -- & -- \\
    \bottomrule
  \end{tabular}
\end{table}

\subsection{Análisis}
\introducirtexto

% ============================================================
\section{Conclusiones}
\label{sec:conclusiones}
% ============================================================

\paragraph{Fuerza bruta.} Su principal ventaja es la simplicidad: es fácil de
implementar, de entender y de demostrar correcta, pues siempre encuentra la
solución óptima. Por ello resulta muy útil como referencia para validar las
demás estrategias. Su gran desventaja es la complejidad exponencial
$\Theta(n \cdot 2^n)$, tanto en tiempo como en espacio en nuestra
implementación, que la hace impracticable para redes con más de unas pocas
decenas de agentes.

\paragraph{Algoritmo voraz.}
\introducirtexto

\paragraph{Programación dinámica.}
\introducirtexto

\paragraph{Comparación general.}
\introducirtexto

% ------------------------------------------------------------
% Referencias
% ------------------------------------------------------------
\begin{thebibliography}{9}
  \bibitem{cormen}
    T.~H. Cormen, C.~E. Leiserson, R.~L. Rivest y C.~Stein.
    \emph{Introduction to Algorithms}, 4.ª ed. MIT Press, 2022.
  \bibitem{enunciado}
    J.~A. Aranda y S.~Galindo.
    \emph{Proyecto de curso: Moderando el extremismo de opiniones en una red social}.
    Análisis y Diseño de Algoritmos II, Universidad del Valle, 2026.
\end{thebibliography}

\end{document}
